\documentclass[a4paper,10pt]{article}

\usepackage[utf8]{inputenc}
\usepackage{t1enc}

\usepackage[utf8]{inputenc}
\usepackage{t1enc}
\usepackage[spanish]{babel}
\usepackage[pdftex,usenames,dvipsnames]{color}
\usepackage[pdftex]{graphicx}
\usepackage{enumerate}
\usepackage{amsmath}
\usepackage{amsfonts}
\usepackage{amssymb}
\usepackage[table]{xcolor}
\usepackage[small,bf]{caption}
\usepackage{float}
\usepackage{subfig}
\usepackage{listings}
\usepackage{bm}
\usepackage{times}

\begin{document}
\setcounter{secnumdepth}{5}
\setcounter{tocdepth}{5}

\renewcommand{\lstlistingname}{C\'odigo Fuente}
\lstloadlanguages{Octave} 
\lstdefinelanguage{MyOctave}[]{Octave}{
        deletekeywords={beta,det},
        morekeywords={repmat}
} 
\lstset{
        language=MyOctave,
        stringstyle=\ttfamily,
        showstringspaces = false,
        basicstyle=\footnotesize\ttfamily,
        commentstyle=\color{gray},
        keywordstyle=\bfseries,
        numbers=left,
        numberstyle=\ttfamily\footnotesize,
        stepnumber=1,                   
        framexleftmargin=0.20cm,
        numbersep=0.37cm,              
        backgroundcolor=\color{white},
        showspaces=false,
        showtabs=false,
        frame=l,
        tabsize=4,
        captionpos=b,               
        breaklines=true,             
        breakatwhitespace=false,      
        mathescape=true
}

%%%%%%%%%%%%%%%%%%%%%%%%%%%%%%%%%%
%%%%%%%% begin TITLE PAGE %%%%%%%%
%%%%%%%%%%%%%%%%%%%%%%%%%%%%%%%%%%
\begin{titlepage}
        \vfill
        \thispagestyle{empty}
        \begin{center}
                \includegraphics{./images/itba_logo.png}
                \vfill
                \Huge{Sistemas de Inteligencia Artificial}\\
                \vspace{1cm}
                \Huge{Métodos de búsqueda no informados y e informados}\\
                \vspace{1cm}
                \Huge{Trabajo Pr\'actico Especial 1}\\
        \end{center}
        \vfill
        \large{
        \begin{tabular}{lcr}
                Civile, Juan Pablo && 50453\\
                Ordano, Esteban && 50753\\
                Crespo, Alvaro && 50758 \\
        \end{tabular}
}
        \vspace{2cm}
        \begin{center}
                \large{25 de marzo de 2013}\\
        \end{center}
\end{titlepage}
\newpage

\setcounter{page}{1}

% \tableofcontents
% \newpage

\section{Definición del problema}

    El problema a resolver es el juego Mahjong (también conocido como Taipei). En particular, este trabajo se focaliza en resolver el juego en su versión de dos dimensiones. \\

    El problema está definido de la siguiente manera:

    \begin{itemize}
        \item Estado inicial: tablero rectangular, representable por una matriz, donde en cada posición hay una ficha de determinado tipo.
        \item Conjunto de acciones posibles: sacar un par de fichas del mismo tipo, que se encuentren en los bordes izquierdo o derecho del tablero.
        \item Modelo de transición: al sacar un par de fichas, se llega a otro estado con un tablero que posee las mismas fichas en las mismas posiciones, excepto por el par de fichas removidas.
        \item Condición de solución: tablero vacío, con todas las fichas removidas.
        \item Función costo de ruta: cantidad de fichas removidas.
    \end{itemize}

\section{Análisis}

    \label{sec:problem_properties}

    De la función de costo y el conjunto de acciones posibles se deduce que toda transición de un estado a otro tiene costo 2.

    Debido a la naturaleza del problema, dado el estado inicial se conoce el costo de ruta de la solución (si es que el tablero es resolvible): el costo será igual a la cantidad de fichas en el tablero.

    Otra característica que presenta el problema es que las soluciones son indistintas unas de otras, ya que todas tendrán el mismo costo de ruta. Esto lleva a que no haya una solución más eficiente que otra.

    Dado un estado con dos acciones posibles, tomar una de ellas conduce a otro estado que siempre contará con la acción que no se tomó en su conjunto de acciones posibles. Como consecuencia, siempre que un estado tenga más de una acción posible (lo cual es muy frecuente, o alcanzar la solución se convierte en un problema trivial), se tendrá más de un camino posible para llegar a una solución, por lo que evitar considerar estados más de una vez es necesario. Dadas $n$ acciones posibles, tomar todas ellas en distinto orden genera $n!$ posibles caminos hasta un estado sin esas $n$ acciones, por lo que la detección de estados equivalentes es considerado indispensable.

\subsection{Heurísticas}

    En esta sección se utilizarán las siguientes definiciones:

    \begin{itemize}
        \item $F_n$: conjunto de fichas en el tablero $n$
        \item $R_n$: cantidad de fichas elegibles para ser eliminadas en el tablero $n$.
        \item $K_n$: conjunto de los distintos tipos de fichas en el tablero $n$.
    \end{itemize}

    \subsubsection{Heurística trivial}

    \[ h(n) = \# F_n \]

    Dada la naturaleza del problema, se conoce exactamente el costo de la solución del problema y no es necesario aproximar este costo. Por esto se denomina a esta heurística como ``trivial''.

    Sin embargo, con el propósito de optimizar \textit{la búsqueda} expandiendo la menor cantidad de nodos, se plantean otras heurísticas cuyo objetivo es expandir aquellos nodos que representan estados considerados con mayor potencial de encontrar una solución.

    \subsubsection{Primer heurística - Ramas angostas}
        
    \[ h(n) = \min \left ( \# F_n, \prod_{i=1}^{\# K_n} {\# c_i \choose 2} \right )\]

    Se propone la productoria de la cantidad de formas de sacar las fichas removibles en el estado para cada tipo de ficha. $c_i$ es el conjunto de fichas de tipo $i$ en los bordes del tablero $n$.

    El objetivo de la heurística es explorar aquellas ramas del árbol de búsqueda donde menos decisiones se deben tomar, evitando la explosión combinatoria (de ahí el nombre "ramas angostas"). A modo de ilustración, si hay tres fichas del mismo tipo en los bordes del tablero, habrá tres posibles elecciones de parejas, y el valor de esta heurística para tal estado será mayor (no será preferido) respecto a un estado en el cual sólo hay dos fichas del mismo tipo en los bordes.

    Para que la heurística sea admisible, se toma el minimo entre el valor de la heurística trivial (que siempre es admisible) y la productoria.
    
    \subsubsection{Segunda heurística - No eliminables}

    \[ h(n) = \# F_n - \# R_n\] 

    Como segunda heurística se propone la resta entre la cantidad total de fichas restantes en el tablero y la cantidad de fichas con posibilidad de ser retiradas en el siguiente movimiento.

    Por construcción esta heurística es admisible, ya que su valor es siempre menor que $\#F_n$, la heurística trivial.

\section{Resolución y resultados}

\subsection{Búsquedas no informadas}

    \subsubsection{Depth First Search}
    \label{sec:dfs}
    Dadas las características del problema, este es el método de búsqueda no informado que mejor se comporta, ya que busca en profundidad y rápidamente llega a estados lejanos al estado inicial, los más cercanos a las soluciones (recurdar que todas las soluciones tienen el mismo costo).

    Un estado del tablero puede ser irresoluble si la última pareja de cada uno de dos o más tipos de fichas se encuentran intercaladas. A modo de ejemplo, el tablero irresoluble de menor cantidad de fichas está formado de una sola línea de forma $ABAB$, con $A$ y $B$ dos tipos de fichas distintos. Nótese que un estado puede ser irresoluble y al mismo tiempo tener un conjunto de acciones posibles no vacío. Cuando esta situación se da, el \textit{DFS} expande una gran cantidad de estados que son irresolubles.

    En el cuadro \ref{averages} se ve que esta estrategia presenta en promedio un buen tiempo de ejecución y una cantidad pequeña de nodos expandidos.
    
    \subsubsection{Breadth First Search}
    En contraposición con el \textit{DFS}, esta estrategia tiene una perfomance pobre, ya que explora muchos estados de más, al hacer una búsqueda ``a lo ancho''. Esta estrategia no es adecuada ya que el problema no presenta una solución óptima o camino óptimo (como ya se observó anteriormente sobre la distancia a la solución), por lo cual analizar todos los caminos posibles no es necesario.

    En los cuadros \ref{crappy-best} y \ref{averages} se observa que las estrategias \textit{DFS} y \textit{Greedy} son significativamente mejores que \textit{BFS}. La primer tabla muestra los mejores resultados obtenidos en 100 ejecuciones con tableros resolvibles generados aleatoriamente. Nótese que incluso cuando estos tableros son fácilmente resueltos por la estrategia \textit{BFS}, \textit{DFS} obtiene mejores resultados.

    \subsubsection{Iterative Deepening}

    Como su definición lo dice, esta estrategia mezcla las características del \textit{DFS} con las del \textit{BFS}, por lo que para este problema, tiene aspectos positivos (los del \textit{DFS}) y también negativos (los del \textit{BFS}). Sin embargo, como consecuencia de las características que comparte con el \textit{BFS}, no es una buena estrategia para resolver el problema.

\subsection{Búsquedas informadas}
    
    \subsubsection{A*}
    
    La naturaleza de las soluciones hace que la estrategia \textit{A*} no sea eficiente, ya que al igual que \textit{BFS}, intenta encontrar la solución de menor costo. Como todas las soluciones tienen el mismo costo, esto a veces implica que con una heurística que no se adapte a esto el algoritmo volverá sobre sus pasos para analizar un estado de menor costo, lo cual pospone el análisis de caminos más cercanos al nivel donde están las soluciones innecesariamente.

    En los cuadros \ref{crappy-best} y \ref{averages} se observa que las estrategias \textit{DFS} y \textit{Greedy} mejoran ampliamente los resultados de \textit{A*}.

    \subsubsection{Greedy Search}
    
    Esta estrategia, por definición, se centra en la heurística elegida y varía su perfomance de acuerdo a ella. A diferencia del \textit{A*}, al considerar el valor de la función de costo del estado actual no hace backtracking innecesario y con una buena heurística, es una estrategia que se adecúa bien al problema.

    El siguiente cuadro muestra una comparación entre 100 ejecuciones de esta estrategia y \textit{DFS}. La variable que comparamos para medir la performance del algoritmo es la cantidad de nodos expandidos.

    La columna \textbf{Peor} cuenta la cantidad de casos que el \textit{Greedy} con esa heurística dio un peor resultado que \textit{DFS}. De la misma manera, \textbf{Igual} y \textbf{Mejor} muestran respectivamente cuantas veces expandieron la misma cantidad de nodos y cuántas veces el \textit{Greedy} expandió menos nodos.

    \begin{table}[H]
        \label{greedy-dfs}

        \begin{center}
        \begin{tabular}{l|r|r|r}
            Heurística & Peor & Igual & Mejor \\
            \hline
            Ramas angostas & 49 & 36 & 15 \\
            No eliminables & 0 & 70 & 30 \\
        \end{tabular}
        \end{center}
        \caption{Comparación entre \textit{Greedy} y \textit{DFS}}

    \end{table}

    Se puede ver en esta tabla que \textit{Greedy - No eliminables} mejora los resultados de \textit{DFS}.
    Una posible explicación para esto es que el uso de la heurística le permite evitar los estados irresolubles (vease la sección \ref{sec:dfs}).

    Si el cálculo de una función que responda si un tablero es irresoluble fuera eficiente en tiempo y memoria, tal función sería una excelente heurística, conduciendo siempre por un camino que lleve a una solución en la menor cantidad de expansiones si es planteada de la forma: $h(n_{irresoluble}) = \infty$, $h(n_{no\ irresoluble}) = \#F_n$.

\section{Conclusiones}
    Dado que se conoce desde un principio la altura del árbol de búsqueda en el cual estarán los nodos que representan estados soluciones (y el costo de transición es constante), interesa que la estrategia de búsqueda llegue rápidamente a esa altura.

    Además, como se sabe que son indistintas (no hay soluciones mejores que otras) basta con hallar una de ellas. Esto lleva a que las estrategias que no intentan minimizar el costo se adecúen mejor al problema.

    Dentro de las estrategias adecuados al problema (\textit{DFS} y \textit{Greedy}), los resultados dependen de la heurística utilizada en el segundo. El problema de ambas estrategias reside en su comportamiento cuando se ingresa a ramas sin solución.  Si la heurística usada en la estrategia \textit{Greedy} ayuda a salir de estas ramas, sus resultados mejoran al \textit{DFS}.

\section*{Anexo: Resultados}

\begin{table}[H]

\begin{tabular}{l|r|r|r|r}
Algoritmo & Tiempo (segs) & Estados & Expandidos & Frontera \\
\hline
DFS & 0.00301599502563 & 31 & 17 & 14 \\
BFS & 0.124269962311 & 472 & 472 & 0 \\
Iterative Deepening 3 & 0.125840902328 & 472 & 472 & 0 \\
Iterative Deepening 5 & 0.128756999969 & 472 & 472 & 0 \\
Greedy Ramas angostas & 0.00515198707581 & 31 & 17 & 14 \\
Greedy No eliminables & 0.00651216506958 & 63 & 17 & 46 \\
A* Ramas angostas & 0.135581016541 & 472 & 471 & 1 \\
A* No eliminables & 0.132388830185 & 472 & 472 & 0 \\

\end{tabular}
\caption{Mejor resultado del \textit{BFS} y \textit{A*}}
\label{crappy-best}

\end{table}

\begin{table}[H]

\begin{tabular}{l|r|r|r|r}
Algoritmo & Tiempo (segs) & Estados & Expandidos & Frontera \\
\hline
A* No eliminables & 21.61545286 & 100443.51 & 27311.39 & 73132.12 \\
A* Ramas angostas & 39.68114223 & 98606.65 & 27098.05 & 71508.6 \\
BFS & 20.0964277 & 100443.51 & 27361.04 & 73082.47 \\
DFS & 0.010342038 & 81.33 & 24.87 & 56.46 \\
Greedy Ramas angostas & 0.018111382 & 128.5 & 21 & 107.5 \\
Greedy No eliminables & 0.02189836 & 130.69 & 36.62 & 94.07 \\
Iterative 3 & 20.72738798 & 99684.76 & 26970.21 & 72714.55 \\
Iterative 5 & 21.6782024 & 100443.51 & 27361.04 & 73082.47
\end{tabular}

\caption{Promedio de los resultados de 100 ejecuciones}
\label{averages}

\end{table}



\end{document}
